\section{Array Programs in C: the MonoCera Tests}\label{sec:monocera}

The MonoCera evaluation (Amilon, Esen, Gurov, Lidström, Rümmer, \emph{Automatic Program Instrumentation for
Automatic Verification}, CAV 2023) has 125 C programs derived from SV-COMP, each with a loop that fills an array and
a loop that checks an assertion over it, one test per aggregation that the assertion admits ($\forall$, $\max$,
$\min$, $\sum$). Unlike the Dafny programs of the previous sections, they carry no specification and no lemmas: the
assertion is the postcondition, and everything a proof needs must come from the database. They are run with all
modules of the database and mathematical integers, as in that evaluation:
\texttt{-unbounded -teq} with the modules, \texttt{-prove -split -prove -split -prove}. The second and third rounds
split the disjunctions that the first introduces, e.g.\ the case of no iteration.

\begin{table}[t]
\centering\small
\begin{tabular}{@{}lrrrrr@{}}
\toprule
aggregation & tests & lookup (\texttt{-prove}) & lookup and Z3 & MonoCera & Horn clauses for Z3 \\
\midrule
$\forall$ & 93 & 53 & 59 & 26 & 0 \\
$\max$ & 4 & 2 & 3 & 2 & 0 \\
$\min$ & 11 & 7 & 9 & 3 & 0 \\
$\sum$ & 17 & 17 & 17 & 10 & 0 \\
\midrule
total & 125 & 79 & 88 & 41 & 0 \\
\bottomrule
\end{tabular}
\caption{Proved tests of the MonoCera evaluation, with mathematical integers: lookup by the database alone, lookup with
the remaining goals lowered for Z3 (\texttt{-fallback}), MonoCera, the instrumenting verifier of that evaluation, and
Horn clauses for Z3 from the same frontend. No result is wrong. The lemmas for these tests started in a separate file
(51 tests in a first run, 73 with that file); their generic part has since moved into the database, and the
specific lemmas (strides, stores ahead) are in \texttt{benchmarks/monocera/lemmas.dfy}, which these runs do not
use.}\label{tab:monocera}
\end{table}


Two things distinguish the C programs. A \texttt{for} loop with a body of conditionals gives one case per path, also
where the paths agree, e.g.\ the counter $i$ with two cases $i + 1$, which the database merges first
(\lemma{trace\_merge\_cases\_ne}). And the harness branches, \texttt{if (N <= 0) return 1}, so the goal has one
branch per path, of which one carries the loops. The iteration constraint of a loop with guarded cases also states
that one of the guards held, e.g.\ $i \bmod 2 = 0 \vee \neg(i \bmod 2 = 0)$: the completeness of the guards in the
context of the update (Section~\ref{sec:engine:cases}) is a fact of the goal.

\subsection{Loops in sequence: sums of filled arrays}\label{sec:brs2}
\begin{lstlisting}[language=C]
for (i = 0; i < N; i++) { if (i % 2 == 0) a[i] = 2; else a[i] = 0; }   // brs2
for (i = 0; i < N; i++) { if (i == 0) sum[0] = 0; else sum[0] = sum[0] + a[i]; }
__VERIFIER_assert(sum[0] <= 2 * N);
\end{lstlisting}
The second loop reads the final array of the first, and the sum is kept in the cell $0$ of a one-element array, so
its trace is a trace of stores at the constant key $0$. The proof needs one lemma per shape, each landing in the
normal form that the next one expects:
\begin{enumerate}[nosep]
\item \emph{The counters.} The traces of $i$ are counters, the constraints and exits give $\intf(n) = N$ for both
  loops, and inside each iteration $i$ is the iteration $k$.
\item \emph{The first loop is stores of a map.} Its update stores a constant chosen by a condition on the index:
\begin{align*}
  &\lemma{stores\_index\_cases\_def}:\\
  &\quad \mathsf{trace}(a_0,\ \mathsf{for}\ \bl{\cnt(0, 1, n)}, \bl{\cnt(0, 1, n)}\\
  &\qquad\quad \mathsf{do}\ a, k, i \mapsto \upd{P(k)}{a[k := c]}\ \ \upd{\neg P(k)}{a[k := d]})\\
  &\quad = \mathsf{stores}(a_0, \bl{\cnt(0, 1, n)}, \mapc{\bl{\cnt(0, 1, n)}}{k \mapsto \ite{P(k)}{c}{d}})
\end{align*}
  with $P$ the parity, abstracted from the guard.
\item \emph{The read in the second loop.} $a'[k]$ is the stored value at $k$ (\lemma{current\_stores}, with the
  membership and position of $k$ in the counter, \lemma{occurs\_counter\_up} and \lemma{indexof\_counter\_up}), decided
  in the context of the iteration, where $0 \le k < \intf(n)$.
\item \emph{The cell is a sum.} The trace of stores into the cell is read as the scalar sum:
\begin{align*}
  &\lemma{select\_current\_trace\_reset\_add}:\\
  &\quad \cur(\mathsf{trace}(a_0, \ldots\ \upd{k = 0}{a[j := c]}\ \ \upd{0 < k}{a[j := a[j] + G(k)]}))[j]\\
  &\quad = \ite{n = 0}{a_0[j]}{c + \mathsf{sum}(\mapc{\bl{\cnt(0, 1, n)}}{k \mapsto \ite{0 < k}{G(k)}{0}})}
\end{align*}
\item \emph{The bound.} \lemma{sum\_cases\_below} bounds a sum of terms that are each $0$, $c$ or $d$ by $e$; its
  left-hand side is the comparison $\mathsf{sum}(\ldots) \le e$, whose right side binds $e$, so the bound need not be
  a fact.
\end{enumerate}
The test \texttt{sum/ms2}, which stores $i \bmod 2$, differs only in the bound (\lemma{sum\_mod\_below}), and
\texttt{sum/conda} only in the first loops: an array filled with ones and then updated in place by cases on the
old value,
\begin{align*}
  &\lemma{stores\_cases\_def}:\\
  &\quad \tr{a_0}{\ldots}{a, k, i \mapsto \upd{P(a[k])}{a[k := c]}\ \ \upd{\neg P(a[k])}{a[k := G(a[k])]}}\\
  &\quad = \mathsf{stores}(a_0, \ldots, \mapc{\ldots}{k \mapsto \ite{P(a_0[k])}{c}{G(a_0[k])}}),
\end{align*}
since each position is read only in its own iteration, before it is written. Its first case stores a constant, not
$G'(a[k])$: by the time the lemma is tried, the read $a[k]$ inside the update has already become $1$ (the stores of
the previous loop, read in the context of the iteration), so the program's $a[k] + 1$ is $2$. A higher-order
instance must contain the term it abstracts, and the constant $2$ contains no $a[k]$.

\paragraph{What this needs.} Loops in sequence compose through their current values: each loop reads the previous
array, the read is resolved in the context of the iteration, and only then is the shape of the next loop recognized.
The lemmas are written for the goal as the engine has simplified it, not for the program text. And where the data has
a different shape, a cell of an array instead of a scalar, the lemma for the new shape lands in the normal form of the
old one (the right-hand side of the scalar sum), so that everything after it is shared. One detail of the engine is
needed here: the instance $(z) \mapsto a[z]$ of a higher-order variable must not be eta-reduced to the array $a$,
since $a(k)$ applied as a function is not the read $a[k]$ that the rules are about; eta reduction applies only to
functions of array type.

\subsection{Strided writes: a quotient in the bound (open)}\label{sec:pr2}
\begin{lstlisting}[language=C]
if (CELLCOUNT % 2 != 0) return 1;   // pr2
for (i = 1; i <= CELLCOUNT / 2; i++) {
  volArray[i*2 - 2] = (3 >= MINVAL) ? 3 : 0;
  volArray[i*2 - 1] = (1 >= MINVAL) ? 1 : 0;
}
for (i = 0; i < CELLCOUNT; i++)
  __VERIFIER_assert(volArray[i] >= MINVAL || volArray[i] == 0);
\end{lstlisting}
The loop stores pairs at $2k$ and $2k + 1$, with values chosen by three cases on \texttt{MINVAL} that do not depend
on the iteration. The module for strided stores reads a cell of the result:
\begin{align*}
  &\lemma{select\_current\_trace\_stride2\_cases}:\quad R \Leftrightarrow (\neg P \wedge \neg Q) \Longrightarrow\\
  &\quad \cur(\mathsf{trace}(a_0, \ldots\ \upd{P}{a[2k := X_1][2k + 1 := Y_1]}\ \ \upd{Q}{\ldots}\ \ \upd{R}{\ldots}))[j]\\
  &\quad = \ite{0 \le j < 2\,\intf(n)}{\ldots}{a_0[j]},
\end{align*}
by the parity of $j$ and the guards. Its side condition requires the guards to be exhaustive; the lemma fixes two
stores and three cases on its left-hand side. It fires on the reads of the checking loop, and every value it
produces satisfies the assertion. The goal that remains is the range above the stores:
\[ 2\,(C / 2) < C \ \wedge\ C \bmod 2 = 0 \ \wedge\ 2\,(C / 2) = i \Longrightarrow \ldots a[i] \ldots \]
for $C$ the size: the loop ran $C / 2$ times, the read at $i = 2\,(C / 2)$ lies beyond the stores, and the case is
impossible since $C$ is even. The arithmetic knows the bounds of $C / 2$, but not that $C \bmod 2$ is
$C - 2\,(C / 2)$, so it does not see the contradiction.

\paragraph{What this needs.} Remainders as linear terms over quotients: with $t \bmod k = t - k \cdot (t / k)$ for a
positive constant $k$, the facts $C \bmod 2 = 0$ and $2\,(C/2) < C$ are inconsistent, and the test closes. A lemma
stating the fact for one form of the equation, e.g.\ $x \bmod 2 = 0 \Rightarrow (n = x / 2 \Leftrightarrow x = 2n)$,
depends on the form in which the engine keeps the equation, which changes with every improvement of the arithmetic;
the remainder belongs to the decision procedure. The tests with $s$ per-cell guards, \texttt{pr3} to \texttt{pr5},
need $2^s$ combinations of cases, which a lemma with a fixed number of cases does not provide: a lemma schema with
the number of stores and cases as a parameter.

\subsection{Stores ahead, in both directions}\label{sec:rewnif}
\begin{lstlisting}[language=C]
for (i = 0; i < SIZE; i++) { if (i + 1 < SIZE) a[i+1] = i; a[i] = i; }       // rewnif
for (i = SIZE - 1; i >= 0; i--) { if (i - 1 >= 0) a[i-1] = i - 2; a[i] = i; } // rewnifrev
for (i = 0; i < SIZE; i++) __VERIFIER_assert(a[i] >= i);
\end{lstlisting}
Each iteration writes one cell ahead before it writes its own; the value written ahead is overwritten by the next
iteration, so a cell holds what its own iteration wrote last, except the cell past the last iteration, which holds
the value written ahead. The module for strided stores has one lemma per direction; for the loop down from $s$,
\begin{lstlisting}
lemma {:rw} {:teq} select_current_trace_ahead_down(a0: array<int, int>, n: Nat, s: int,
                                                   F: (int) -> int, G: (int) -> int, j: int)
  ensures
    current(trace(a0, for ⌊counter(0, 1, n)⌋, ⌊counter(s, -1, n)⌋ do a, k, i ::
              case 1 <= i => a[i + (-1) := F(i)][i := G(i)]
              case i < 1 => a[i := G(i)]))[j]
      == if s - integer(n) < j && j <= s then G(j)
         else if j == s - integer(n) && 0 <= j && 0 < integer(n) then F(j + 1)
         else a0[j]
\end{lstlisting}
It fires on the reads of the checking loop with $s := \texttt{SIZE} - 1$, $F := (z) \mapsto z - 2$ and
$G := (z) \mapsto z$, and the assertion holds in every branch; the test closes under $\forall$ and $\min$.

\paragraph{What this needs.} The pattern contains two literals that a matcher must treat modulo the normal form of the
goal: the step $-1$ of the counter dependency, matched through the definition of the program's counter, and the store
position $i + (-1)$, which the goal has as $i - 1$, a sum with a negative literal. The engine compares a negative
literal among the summands as the negation of its absolute value, so the lemma can be written as the program reads.
Without that, each such position needs a higher-order variable $P(i)$ pinned to $i - 1$ by a side condition, and each
literal step a variable $d$ with $d = -1$. The third variant of the test, \texttt{rewnifrev2}, writes ahead upwards in a
loop that runs downwards, a shape for which there is no lemma.

\subsection{Recurrences on the previous cell}\label{sec:ifeqn2}
\begin{lstlisting}[language=C]
for (i = 0; i < N; i++) { if (i == 0) a[0] = 4; else a[i] = a[i-1] + 4; }   // ifeqn2
for (i = 0; i < N; i++) { if (i == 0) b[0] = 1; else b[i] = b[i-1] + a[i-1]; }
for (i = 0; i < N; i++) __VERIFIER_assert(b[i] == 2*i*i + 2*i + 1);
\end{lstlisting}
A store at $k$ of a value read at $k - 1$ is not a store of a map: the value depends on the trace, not on the initial
array. The aggregation module names the recurrence by its closed form,
\begin{lstlisting}
lemma {:rw} {:teq} select_current_trace_recurrence(a0: array<int, int>, n: Nat, c: int, j: int,
                                                   G: (int) -> int)
  ensures current(trace(a0, for ⌊counter(0, 1, n)⌋ do a, i ::
            case i == 0 => a[0 := c] case 0 < i => a[i := a[i - 1] + G(i)]))[j]
    == if 0 <= j && j < integer(n)
       then c + sum(maps(for ⌊counter(1, 1, natural(j))⌋ yield k :: G(k)))
       else a0[j]
\end{lstlisting}
The first loop gives $a'[j] = 4 + \mathsf{sum}(\mapc{\bl{\cnt(1, 1, \natf(j))}}{k \mapsto 4}) = 4 + 4j$ (the map of a
constant and its sum). In the second loop the summand $G(k) = a'[k - 1]$ is then $4k$, so
$b'[j] = 1 + \mathsf{sum}(\mapc{\bl{\cnt(1, 1, \natf(j))}}{k \mapsto 4k})$; the map of a multiple of the counter is
a counter, $\bl{\cnt(4, 4, \natf(j))}$ (\lemma{maps\_counter\_scale}), whose sum is the arithmetic series
(\lemma{sum\_counter}), $2j^2 + 2j$. The read of the checking loop at its index meets the assertion exactly.

\paragraph{What this needs.} The read $a[i - 1]$ of the pattern must match the goal's $a[k + (-1)]$, the same negative
literal as above; the higher-order summand $G$ must take the read of the other array as a whole; and the closed forms
of maps and sums over counters must reach polynomials, which the arithmetic compares. The recurrences with other
increments, products or conditions (\texttt{ifeqn3} to \texttt{ifeqn5}, \texttt{sqm}, \texttt{condm}) need their own
closed forms; a recurrence solver as a generator of such lemmas would cover the polynomial ones.
